\documentclass{beamer}

% Uso de estilos en beamer
%\usetheme{Frankfurt}
%\usetheme{JuanLesPins}
\usetheme{Darmstadt}
 %AnnArbor
 %Antibes
 %Berlin
 %Boadilla
 %Copenhagen
 %Darmstadt
 %default
 %Dresden
 %Frankfurt
 %Hannover
 %Ilmenau
 %JuanLesPins
 %Madrid
 %Rochester
 %Warsaw
 %}
\useoutertheme{}
\useinnertheme{}
\usefonttheme{}
\usecolortheme{}
%beaver
%beetle
%crane
%default
%dolphin
%dove
%fly
%seagull
%sidebartab

%\setbeamertemplate{background canvas}{\includegraphics[width=\paperwidth,height=\paperheight]{pix/bg}}

% Desactivar la barra de navegacion
%\setbeamertemplate{navigation symbols}{}

\usepackage[spanish]{babel}
\usepackage[utf8]{inputenc}

% Redefinicion de teorema, definicion, ejemplo
\newtheorem{teorema}[theorem]{Teorema}
\newtheorem{definicion}[theorem]{Definición}
\newenvironment{ejemplo}{\begin{exampleblock}{Ejemplo}}{\end{exampleblock}}

\title{Maestro Virtual}
\subtitle{Novela gráfica creada con el motor Ren'Py}
\author[Milton Mazzarri]{Milton Mazzarri \\
\href{mailto:milmazz@gmail.com}{milmazz@gmail.com} \\
}
\institute[gUsLA]{Grupo de Usuarios de Software Libre de la Universidad de Los Andes \\
%
\date{11 de Julio del 2008}

\hypersetup{colorlinks=true,linkcolor=blue}

\hypersetup{
pdfauthor={Milton R. Mazzarri S.},
pdftitle={Maestro Virtual: Novela gráfica creada con el motor Ren'Py}
}

\begin{document}

\begin{frame}
\titlepage
\end{frame}

\section{Maestro Virtual}

\subsection{Origen}
\frame{
  \frametitle{El origen del proyecto}

  \begin{itemize}
    \item
      El \alert{Maestro Virtual} nace como colaboración con el proyecto educativo del grupo de investigación en Sistemología Interpretativa de la Universidad de Los Andes.
    \item
      El proyecto educativo está basado en la enseñanza a los niños a través de narraciones.
    \item
      La estrategia es incluir los contenidos en el mismo cuento, aprovechando la narración de sucesos para incluir actividades de enseñanza. 
    \item
      El propósito es hacer uso del conocimiento para enriquecer la historia.
  \end{itemize}
}

\subsection{El contenido}
\frame{
  \frametitle{El contenido de los cuentos}

  \begin{itemize}
    \item
      La Dra. Miriam Villarreal (ULA) diseñó los contenidos que dicen paso a paso como realizar las actividades en torno al cuento. 
    \item
      En este caso, \emph{El Príncipe Feliz} de Oscar Wilde.
    \item
      El diseño del cuento de \emph{El Príncipe Feliz} permitió crear una guía para niños de quinto grado de primaria.
  \end{itemize}
}

\subsection{Uso de Tecnología de Información}
\frame{
  \frametitle{Inclusión de tecnologías de información}

  \begin{itemize}
    \item
      El proyecto \alert{Maestro Virtual} pretende desarrollar una aplicación de \emph{software} cuyo guión sea el contenido para quinto grado diseñado sobre el cuento \emph{El Príncipe feliz}.
    \item
      Para lograr esto se comenzó a observar y estudiar los principios educativos en torno al proyecto educativo, seguido de una investigación acerca de las herramientas \alert{libres} existentes para poder desarrollar la aplicación de manera que exista una interacción entre el niño y la máquina.
    \item
      El \emph{software} \alert{Maestro Virtual} no pretende sustituir al maestro, será un apoyo. El hecho de usar el computador como herramienta atraerá la atención del niño.
  \end{itemize}
}

% TODO
%\section{Novelas gráficas}

\section{Ren'Py}

\subsection{Engine de Ren'Py}
\frame{
  \frametitle{Ren'Py: ¿Qué es exactamente?}

  \begin{itemize}
    \item
      Es un motor que soporta la creación de novelas gráficas, simuladores de citas y distintas formas de contar historias a través de un ordenador.
    \item
      El motor permite crear novelas personalizables y flexibles.
    \item
      Ofrece opciones de cualquier novela gráfica.
      \begin{itemize}
        \item
          Guardar, cargar partidas.
        \item
          Preferencias.
        \item
          Retroceder escenas.
        \end{itemize}
    \item
      La licencia de Ren'Py permite crear novelas comerciales y no-comerciales.
    \item
      Ren'Py funciona en Linux, Mac y \emph{Winbugs}.
  \end{itemize}
}

\subsection{Funcionalidades}
\begin{frame}[fragile]
  \frametitle{Dialogos, Narraciones y Lógica de la novela}

  \begin{itemize}
    \item
      La sintaxis de las novelas gráficas usando Ren'Py son muy similares al formato de un guion de una película.
    \item
      Ren'Py necesita de poco código para que tus personajes realmente digan algo.
  \end{itemize}

  \begin{ejemplo}<3->
\begin{verbatim}
# script.rpy
init:
    $ pebbels = Character("Pebbels", color="#c8ffc8")

label start
    "Hace un año..."
    pebbels "Gazoo"
    "Bamm-Bamm" "Bam Bam..."
\end{verbatim}
  \end{ejemplo}
\end{frame}

\frame{
  \frametitle{Dialogos, Narraciones y Lógica de la novela}
 
  \begin{itemize}
    \item
      En los scripts de Ren'Py los personajes están referenciados por nombre, no por número.
    \item
      Ren'Py ajusta el texto de las narraciones o dialogos para que encaje en la proyección de pantalla.
    \item
      La lógica de la novela puede ser basada en menu de opciones o variables, de ese modo una novela gráfica puede tener N finales, mejorando la experiencia del usuario.
    \item
      La presentación de las narraciones es \alert{totalmente} personalizable, ubicación, colores, estilos, fuentes tipográficas, imágenes.
    \item
      El texto puede incluir etiquetas.
  \end{itemize} 
}

\frame{
  \frametitle{Interfaz de usuario y estado de la novela gráfica}

  \begin{itemize}
    \item
      Soporta tanto guardar/cargar sesiones automáticamente.
    \item
      Proporciona múltiples \alert{slots} para almacenar estados.
    \item
      Interacción con el mouse, teclado y \alert{joystick}.
    \item
      Puede programar el avance automático de las escenas, de ese modo los usuarios no necesitan dar click para avanzar.
  \end{itemize}
}

\frame{
  \frametitle{Efectos especiales}

  \begin{itemize}
    \item
      Soporte de modo pantalla completa o en una ventana.
    \item
      Un número ilimitado de imágenes puede ser mostradas en una escena.
    \item
      Se pueden reproducir videos \texttt{mpeg1}.
    \item
      Soporta la precarga de imágenes, de ese modo se evitan los retardos en la presentación de las escenas.
    \item
      Proporciona efectos en las transiciones de las escenas.
  \end{itemize}
}

\frame{
  \frametitle{Sonidos}

  \begin{itemize}
    \item
      Reproducción de ficheros \texttt{mp3}, ogg vorbis, entre otros.
    \item
      Efectos aparecer/desvanecer progresiva de audio.
    \item
      Soporte para voces (actuación de los dialogos y narraciones).
    \item
      Efectos de sonido, música y voces puede reproducirse simultáneamente.
  \end{itemize}
}

\frame{
  \frametitle{Distribución}

  \begin{itemize}
    \item
      Soporta Linux, Mac y el \emph{innombrable}. Soporte para otras plataformas es posible.
    \item
      Ren'Py es distribuido bajo la licencia MIT, eso quiere decir que puede vender las novelas hechas sin necesidad de pagar por Ren'Py. La distribución libre obviamente es permitida.
    \item
      Los script de Ren'Py puede ser ofuscados para su distribución, haciendo dificil la lectura o modificación del código fuente.
    \item
      Los ficheros de medios pueden ser archivados para la distribución de tal modo que sea dificil acceder a los ficheros individualmente.
  \end{itemize}  
}

\subsection{Comenzando el desarrollo}
\frame{
  \frametitle{Desarrollo de novelas}

  \begin{columns}

  \begin{column}{5cm}
    \begin{itemize}
      \item<1->
        La distribución de Ren'Py viene con un lanzador que le ayuda a comenzar y manejar sus proyectos.
      \item<3->
        Ren'Py le ofrece algunas plantillas para facilitarle el comienzo.
      \item<5->
        La documentación de Ren'Py se encuentra en el sitio oficial.
    \end{itemize}
  \end{column}

  \begin{column}{5cm}
    \begin{overprint}
      \includegraphics<2>[scale=0.5]{pix/launcher}
      \includegraphics<4>[scale=0.5]{pix/template}
      \includegraphics<6>[scale=0.3]{pix/renpy-site}
    \end{overprint}
  \end{column}

  \end{columns}
}

\subsection{¿Cómo encontrar ayuda?}
\frame{
  \frametitle{Ayuda en línea}

  \begin{itemize}[<+->]
      \item
        Los foros de \href{http://lemmasoft.renai.us/forums/viewforum.php?f=8}{Lemmasoft} le pueden ayudar en el desarrollo de las novelas.
      \item
        Ayuda en línea en el canal IRC \alert{\#renpy} en \href{http://freenode.net/}{freenode.net}
      \item
        Revisión del código fuente disponible para las novelas gráficas. Hasta ahora son más de \alert{110} novelas hechas con Ren'Py.
  \end{itemize}
}

\section{Involucrados en el Proyecto Maestro Virtual}
\frame{
  \frametitle{Autores}

  Hasta ahora los involucrados en el Proyecto \alert{Maestro Virtual} son:

  \begin{itemize}
    \item
      José Aguilar -- \href{mailto:jaguilar@cenditel.gob.ve}{jaguilar@cenditel.gob.ve} -- Asesor 
    \item
      Ramsés Fuenmayor -- \href{ramsesfa@ula.ve}{ramsesfa@ula.ve} -- Asesor
    \item
      Ana Rangel -- \href{mailto:arangel@cenditel.gob.ve}{arangel@cenditel.gob.ve} -- Investigación, voces.
    \item
      Luz Mairet Chourio \href{mailto:lchourio@cenditel.gob.ve}{lchourio@cenditel.gob.ve} -- Investigación.
    \item
      Iván Dario Hernández -- \href{mailto:ivandarioh@yahoo.es}{ivandarioh@yahoo.es} -- Arte.
    \item
      Milton Mazzarri -- \href{mailto:milmazz@milmazz.com.ve}{milmazz@milmazz.com.ve} -- Programador, voces.
  \end{itemize}  
}

\section{Herramientas utilizadas en el desarrollo}
\frame{
  \frametitle{Herramientas utilizadas}
  
  \begin{columns}
  
  \begin{column}{5cm}
    \begin{itemize}
      \item<1->
        Ren'Py
      \item<2->
        Gimp
      \item<3->
        Inkscape
      \item<4->
        Audacity
      \item<5->
        Scite
      \item<6->
        Vim
    \end{itemize}
  \end{column}

  \begin{column}{5cm}
    \begin{overprint}
      \includegraphics<2>{pix/wilber}
      \includegraphics<4>{pix/audacity}
      \includegraphics<6>{pix/vim}
    \end{overprint}
  \end{column}
  \end{columns}
}

% TODO: Colocar el portal del svn, trac, lista de distribución, etc.
\end{document}
